\appendix

\chapter{Extended Results and Result Videos}
\label{app:results}

The recordings below accompany this report. They are grouped by what they compare: within each
comparison table one row shows the same scene under two or more pipeline variants, so the effect of
a single change can be seen directly. Scene names describe the captured interaction (for example
``cloth, single lift'' is a cloth lifted from one point). Recording links are given in shortened
form.

\section{Improvement Chain on the Sloth Sequence}
\label{sec:app-vid-sloth-chain}

Table~\ref{tab:app-vid-sloth} shows the improvement chain on the sloth sequence: the original
pipeline (sensor reference depth), the pipeline with DA3 inferred depth, and the pipeline with
smoothed DA3 depth. Each stage replaces one component of the previous one. The measurements in
Chapter~\ref{cha:evaluation} quantify the last step: for this sequence the smoothed depth keeps the
agreement with the sensor reference (Spearman 0.988, ordinal error 3.3\% on camera 0) with a higher
gradient correlation than raw DA3 (0.169 against 0.117; Table~\ref{tab:eval-depth-sensor}).

\begin{table}[H]
    \caption{Sloth sequence: the three pipeline stages.}
    \label{tab:app-vid-sloth}
    \begin{center}
        \small
        \begin{tabular}{@{}p{6.2cm}l@{}} \toprule
            Stage & Recording \\ \midrule
            1. Original pipeline (reference depth) & \url{youtu.be/UyQw7Hx25D8} \\
            2. DA3 inferred depth with streaming   & \url{youtu.be/Wldxhw5DKHs} \\
            3. Smoothed DA3 inferred depth         & \url{youtu.be/xjoEsQoWlXw} \\ \bottomrule
        \end{tabular}
    \end{center}
\end{table}

\section{Effect of the Temporal Depth Smoother}
\label{sec:app-vid-smoother}

\paragraph{Raw DA3 depth against smoothed depth.}
\label{sec:app-vid-raw}

Each row of Table~\ref{tab:app-vid-raw} shows one scene twice: the raw DA3 depth prediction and the
same frames after temporal smoothing. The raw recordings show the frame-to-frame flicker that
motivates the smoother (Chapter~\ref{cha:smoother}): the depth of static surfaces shifts between
consecutive frames and the object--background boundary moves.

\begin{table}[H]
    \caption{Raw DA3 depth against smoothed depth.}
    \label{tab:app-vid-raw}
    \begin{center}
        \small
        \begin{tabular}{@{}p{3.4cm}ll@{}} \toprule
            Scene & Raw DA3 depth & Smoothed depth \\ \midrule
            Cloth, single lift & \url{youtu.be/c_K8Iw4v8GM} & \url{youtu.be/Oix-bPoGo1Y} \\
            Zebra, double lift & \url{youtu.be/22x_s_dwmJA} & \url{youtu.be/jrlRHbfc8wo} \\
            Mat, single lift   & \url{youtu.be/CfmDhliW4e4} & \url{youtu.be/eaA4PYEcdmU} \\ \bottomrule
        \end{tabular}
    \end{center}
\end{table}

\paragraph{Pipeline replay before and after smoothing.}
\label{sec:app-vid-smooth}

Table~\ref{tab:app-vid-smooth} shows the pipeline replay of the same two scenes before and after
the inferred depth was smoothed, at matched times.

\begin{table}[H]
    \caption{Pipeline replay before and after depth smoothing.}
    \label{tab:app-vid-smooth}
    \begin{center}
        \small
        \begin{tabular}{@{}p{3.4cm}ll@{}} \toprule
            Scene & Before smoothing & After smoothing \\ \midrule
            Cloth, single lift  & \url{youtu.be/oMyWfrkThjA} & \url{youtu.be/sFhpqg2umsA} \\
            Sloth, double stretch & \url{youtu.be/Pyi-6zr_SVY} & \url{youtu.be/DJOm86ircG4} \\ \bottomrule
        \end{tabular}
    \end{center}
\end{table}

\section{Depth-Source Variants}
\label{sec:app-vid-depth-variants}

Table~\ref{tab:app-vid-variants} lists further depth sources used for the same inference
comparison. The reference and DA3 rows repeat the first two stages of
Table~\ref{tab:app-vid-sloth} so that the set is self-contained.

\begin{table}[H]
    \caption{Inference with alternative depth sources.}
    \label{tab:app-vid-variants}
    \begin{center}
        \small
        \begin{tabular}{@{}p{4.6cm}ll@{}} \toprule
            Depth source & Recording & Notes \\ \midrule
            Reference (sensor depth) & \url{youtu.be/UyQw7Hx25D8} & original pipeline \\
            DA3 inferred depth with streaming & \url{youtu.be/Wldxhw5DKHs} & adopted depth source \\
            VDA                      & \url{youtu.be/GUog6wd8YiI} & video model oracle \\
            Without streaming        & \url{youtu.be/NbU0TzVzMGM} & without chunk stitching \\ \bottomrule
        \end{tabular}
    \end{center}
\end{table}

\section{Streaming Inference}
\label{sec:app-vid-streaming}

Table~\ref{tab:app-vid-streaming} compares depth inference with and without chunk stitching
(streaming), for the final data and the final matching visualisations, with a reference recording
of the final data.

\begin{table}[H]
    \caption{Streaming inference against non-streaming inference.}
    \label{tab:app-vid-streaming}
    \begin{center}
        \small
        \begin{tabular}{@{}p{3.4cm}lll@{}} \toprule
            Visualisation & Without streaming & With streaming & Reference \\ \midrule
            Final data & \url{youtu.be/wtfQiYJZQfQ} & \url{youtu.be/b3FfEog8buc} & \url{youtu.be/bH6H2vx3Q4k} \\
            Final matching & \url{youtu.be/m4g420lpaH4} & \url{youtu.be/M2VwlY6B9j4} & --- \\ \bottomrule
        \end{tabular}
    \end{center}
\end{table}

\section{Tracking Comparison with SpatialTrackerV2}
\label{sec:app-vid-spatracker}

Table~\ref{tab:app-vid-spatracker} compares the tracked result of the same monkey sequence with the
2D-tracking-based pipeline and with SpatialTrackerV2 \parencite{spatrackerv2}.

\begin{table}[H]
    \caption{Monkey sequence: DIPhysTwin against SpatialTrackerV2.}
    \label{tab:app-vid-spatracker}
    \begin{center}
        \small
        \begin{tabular}{@{}p{3.6cm}ll@{}} \toprule
            Variant & Recording & Notes \\ \midrule
            DIPhysTwin & \url{youtu.be/8WwHEsOErsI} & spatial separation on fast movement \\
            SpatialTrackerV2 & \url{youtu.be/DijnnXjS_Xc} & frequent 3D artefacts \\ \bottomrule
        \end{tabular}
    \end{center}
\end{table}

\section{Environment Interaction: Rope with and without Ground Alignment}
\label{sec:app-vid-ground}

Table~\ref{tab:app-vid-ground} shows the rope sequence before and after the ground-plane alignment
was added (Chapter~\ref{cha:system-design}), together with the reference recording of the same
interaction.

\begin{table}[H]
    \caption{Rope sequence: stiff result, reference interaction, and recovered ground contact.}
    \label{tab:app-vid-ground}
    \begin{center}
        \small
        \begin{tabular}{@{}p{5.2cm}ll@{}} \toprule
            Variant & Recording & Notes \\ \midrule
            Without ground alignment & \url{youtu.be/I59sUpi_Rvk}  & rope overly stiff \\
            Reference interaction    & \url{youtu.be/fhun0051Rps}  & rope rests on the ground \\
            With ground alignment    & \url{youtu.be/zXyEl8aMBYk}  & ground contact recovered \\ \bottomrule
        \end{tabular}
    \end{center}
\end{table}

\section{Individual Demonstrations}
\label{sec:app-vid-solo}

Single recordings without a paired comparison, grouped by type.

\begin{table}[H]
    \caption{Individual demonstration recordings.}
    \label{tab:app-vid-solo}
    \begin{center}
        \small
        \begin{tabular}{@{}llp{4.6cm}@{}} \toprule
            Recording & Scene & Notes \\ \midrule
            \multicolumn{3}{@{}l}{\emph{Tracking results}} \\
            \url{youtu.be/94h6tWvJEwE} & Mat, double bend & tracking reconstruction \\
            \url{youtu.be/ICQz9aUabAo} & Mat, single lift & tracking reconstruction \\ \midrule
            \multicolumn{3}{@{}l}{\emph{Pipeline demonstrations, PhysTwin data}} \\
            \url{youtu.be/nsndv85sxcU} & Sloth, double stretch & \\
            \url{youtu.be/t97kz2DZZxk} & Zebra, double lift & \\
            \url{youtu.be/wVsCuG4sxWI} & Cloth, double hand lift & \\
            \url{youtu.be/FUNXNBtAupQ} & Cloth, single lift & degraded result \\
            \url{youtu.be/I59sUpi_Rvk} & Rope, single lift & see Table~\ref{tab:app-vid-ground} \\ \midrule
            \multicolumn{3}{@{}l}{\emph{Pipeline demonstrations, self-collected data}} \\
            \url{youtu.be/U8ICkg25_4U} & Mat, double bend & \\
            \url{youtu.be/rtjAy8hZE8k} & Mat, single lift & \\
            \url{youtu.be/RTn-0Zg8wAA} & Water bottle, single lift & poor result \\
            \url{youtu.be/SbAI6K0lMDk} & T-shirt, single lift & \\
            \url{youtu.be/svUzZCJdgDk} & Monkey, double bend & \\
            \url{youtu.be/yZ0TzMy9XA0} & Package, single lift & \\ \midrule
            \multicolumn{3}{@{}l}{\emph{New-data captures, three camera corners}} \\
            \url{youtu.be/sFbwJNk7rUs} & Mat, single lift & camera corner 0 \\
            \url{youtu.be/LdueFxEb7H8} & Mat, single lift & camera corner 1 \\
            \url{youtu.be/ZcDswgNRI40} & Mat, single lift & camera corner 2 \\ \midrule
            \multicolumn{3}{@{}l}{\emph{Issue and failure recordings}} \\
            \url{youtu.be/e6UjWpdsnjg} & Mat, double bend & object split during interaction \\
            \url{youtu.be/8WwHEsOErsI} & Monkey, double bend & see Table~\ref{tab:app-vid-spatracker} \\
            \url{youtu.be/4lE1MX_xUnM} & Rope, single lift & depth flicker propagating into 3D tracks \\ \midrule
            \multicolumn{3}{@{}l}{\emph{Further recordings}} \\
            \url{youtu.be/TVVKJ1hBin8} & Sloth, double stretch & interactive playground session \\
            \url{youtu.be/cywRH67oSGY} & Sloth, double stretch & depth-quality comparison with metrics \\
            \url{youtu.be/R-0QuXLqUTQ} & Mat, double bend & early result, poor \\
            \url{youtu.be/44K07VyE1iI} & Mat, single lift & early result \\ \bottomrule
        \end{tabular}
    \end{center}
\end{table}

\chapter{Hyper-parameters and Loss Weights}
\label{app:hyperparameters}

This appendix records the settings needed to reproduce the temporal depth smoother. The values
are the training configurations saved with the checkpoints.

\begin{table}[H]
    \caption{Training and inference hyper-parameters of the reported smoother run.}
    \label{tab:app-hyperparameters}
    \begin{center}
        \begin{tabular}{@{}p{3.6cm}p{10.4cm}@{}}
            \toprule
            Setting & Value \\
            \midrule
            Optimiser & Adam; weight decay $10^{-4}$ \\
            Learning rate & $3\times10^{-3}$ \\
            Schedule & ReduceLROnPlateau; factor $0.5$, patience 5 \\
            Batch size & 2 \\
            Epochs & 50; checkpoints saved for epochs 0--49 \\
            Precision & CUDA autocast with bfloat16 \\
            Hardware & NVIDIA GeForce RTX 3060 \\
            Inference chunk length & 16 frames per chunk, matching the training window; chunks
            overlap and are blended with a weighted window \\
            \bottomrule
        \end{tabular}
    \end{center}
\end{table}

\begin{table}[H]
    \caption{Loss weights of the best saved smoother runs; the first row is the run behind the
    reported results. $\lambda_{\mathrm{edge}}$ and $\lambda_{\mathrm{grad}}$ are the
    edge-sharpness and spatial-gradient terms of Section~\ref{sec:smoother-losses}.}
    \label{tab:app-lossweights}
    \begin{center}
        \small
        \begin{tabular}{@{}p{5.6cm}ccccc@{}}
            \toprule
            Run & $\lambda_{\mathrm{fid}}$ & $\lambda_{\mathrm{tgm}}$ & $\lambda_{\mathrm{tv}}$
            & $\lambda_{\mathrm{edge}}$ & $\lambda_{\mathrm{grad}}$ \\
            \midrule
                        \texttt{best\_model} & 1.2 & 0.5 & 0.5 & 0.3 & 0 \\
            \texttt{best\_model\_50\_epochs} & 1.2 & 0.5 & 0.5 & 0.3 & 0 \\
            \texttt{best\_model\_50\_epochs\_no\_edge\_blur} & 1.3 & 0.8 & 0.8 & 0.05 & 0.04 \\
            \texttt{best\_model\_new\_data} & 1.3 & 0.8 & 0.8 & 0.05 & 0.05 \\
            \bottomrule
        \end{tabular}
    \end{center}
\end{table}

The temporal spans follow the code defaults used by all runs: TGM $k = 1, 2, 3$ with weights
$1, 0.25, 0.07$, and TV $k = 1, 2$ with weights $1, 0.25$.



\chapter{Source Code}
\label{app:code}

The implementation is released at
\url{https://github.com/Hanson-Sun/PhysTwin/tree/depth-inference-phystwin}.

